\PassOptionsToPackage{svgnames}{xcolor}
\documentclass[a4paper,12pt]{article}
\usepackage[utf8]{inputenc}
\usepackage[bulgarian]{babel}
\usepackage{hyperref}
\usepackage{subfiles}
\usepackage{svg}
\usepackage{amsmath}
\usepackage[final]{pdfpages}
\usepackage[table,xcdraw]{xcolor}

\usepackage{longtable}
\usepackage{amsmath}
\usepackage{algorithm}
\usepackage[noend]{algpseudocode}

\usepackage{graphicx}
\graphicspath{ {./images/} }
\makeatletter
\def\BState{\State\hskip-\ALG@thistlm}
\makeatother
\hypersetup{
	colorlinks,
	linkcolor={black},
	citecolor={black},
	urlcolor={black}
}
\usepackage{lmodern}
%\usepackage{geometry}
%\usepackage{layout}
\usepackage{graphicx}
\usepackage{wrapfig}

\title{WebFrame++}
\author{Алекс Цветанов}

\newcommand{\titleGP}{\begingroup
	\centering
	%\vspace{\baselineskip}
	%\begin{wrapfigure}{r}{\textwidth}
    %\end{wrapfigure}
		\includegraphics[width=\textwidth]{images/logo.png}
    \\\vspace{0.1cm}
	\rule{\textwidth}{1.6pt}\vspace{-\baselineskip}\vspace{2pt}
	\rule{\textwidth}{0.4pt}\\[\baselineskip]
	
	{\LARGE WebFrame++ \\ \small{C++ web application framework}}\\[0.2\baselineskip]
	
	\rule{\textwidth}{0.4pt}\vspace{-\baselineskip}\vspace{3.2pt}
	\rule{\textwidth}{1.6pt}\\[\baselineskip]
    \includegraphics[width=0.35\textwidth]{images/webpp-logo.png}

	\vspace{20pt}
	{
	{\LARGE Алекс Цветанов}\\
    {\normalsize Софийска Математическа Гимназия, София, България}
	\par}
	
	\vspace{2\baselineskip}
	
	\vfill
	
	{\scshape Ученическия институт по математика и информатика на БАН, 2020} \\[0.3\baselineskip]
	{\large Боровец, България}\par
	
	\endgroup}


\usepackage{tcolorbox}
\newenvironment{myblock}[1]{%
	\tcolorbox[beamer,%
	noparskip,breakable,
	colback=LightGreen,colframe=DarkGreen,%
	colbacklower=LimeGreen!75!LightGreen,%
	title=#1]}%
{\endtcolorbox}
\newenvironment{myblocksecond}[1]{%
	\tcolorbox[beamer,%
	noparskip,breakable,
	colback=Yellow,colframe=Brown,%
	colbacklower=Brown!75!Yellow,%
	title=#1]}%
{\endtcolorbox}

\begin{document}
	
	\titleGP
	\newpage
	
	\tableofcontents
	\newpage
	
	\begin{abstract}
    \begin{text}
    \par
    WebFrame++ е съвременна библиотека, която цели да направи един от най-бързите програмни езици (C++) удобен за разработка на интернет приложения.\par
    Разработката предоставя удобен синтаксис за създаване на уеб сървър, задаване на пътища и тяхното съдържание по начин, наподобяваш този на повечето използвани в момента библиотеки и модули със същата цел, но за други езици.\par
    Примери за такива са Express.JS за Node.JS, Flask и Django за Python, дори Spring framework за Java.\par
    \vspace{1cm}
    WebFrame++ is a modern library that aims to make
one of the fastest programming languages (C ++) efficient for web applications. \par
The library provides a convenient syntax for creating a web server, setting up paths and the content in a way similar to most currently used libraries and
modules for the same purpose but for other languages.\par
Examples of these are Express.JS for Node.JS, Flask and Django for
Python and even Spring framework for Java.\par



    \end{text}
	\end{abstract}
	\newpage
	
	\section{Въведение}
	    \subfile{chapters/introduction}
	\section{Проблем}
	    \subfile{chapters/problem}
	\section{Решение}
	    \subfile{chapters/solution}
	\section{Сравнение} \label{compare}
	    \subfile{chapters/compare}
	\section{Технологии}
	    \subfile{chapters/techniques}
	\section{Бъдеща разработка и разпространение}
	    \subfile{chapters/future}
	\section{Благодарности}
    	Благодарности за помощта за оформянето на документацията:
    	\begin{itemize}
    		\item Надежда Цачева
    		\item Радостина Богословова
    	\end{itemize}
    	Благодаря още на:
    	\begin{itemize}
    		\item Ученическия институт на БАН
    		\item Българска академия на науките
    	\end{itemize}

    \bibliographystyle{plain}
    \bibliography{ref}
	
\end{document}
